% Authority: exact-scope leaf 17; full PDF page 96; printed p.79.
% U:p079.r28.sec.005 | heading; R28 sec.005 visually replayed
\sectionnumber{5}
% U:p079.r28.par.020 | paragraph-start; R28 par.020 visually replayed
\par\indent
Let us denote by $a$ a prime number $4n+1$, and by $p$ a prime number likewise $4n+1$, and
moreover capable of being put under the form $t^2-au^2$. We can therefore put $p=t^2-au^2$ where
the number $t$, as is easy to see, is necessarily odd. One will find, as in paragraph 3, that if
one denotes by $2^\nu$ the highest power of 2 that divides $u$, one has
$\leg{u}{p}=\leg{2^\nu}{p}$, a relation which one will then change, as in the cited place, into
this one: $\leg{u}{p}=(-1)^{\frac{p-1}{4}}$.
% U:p079.r28.par.021 | paragraph-start; R28 par.021 visually replayed
\par\indent
Let us now decompose the odd number $t$ into its simple factors, putting $t=gg'g''\ldots{}\times
hh'h''\ldots$, where the prime numbers $g$, $g'$, $g''$, etc., $h$, $h'$, $h''$, etc. are such
that:
% U:p079.r28.eq.014 | continuation; R28 eq.014 visually replayed
\begin{equation*}\tag{$\varepsilon$}
\left\{
\begin{array}{llll}
\leg{-a}{g}=1,& \leg{-a}{g'}=1,& \leg{-a}{g''}=1,& \text{etc.}\\[0.3em]
\leg{-a}{h}=-1,& \leg{-a}{h'}=-1,& \leg{-a}{h''}=-1,& \text{etc.}
\end{array}
\right.
\end{equation*}
% U:p079.r28.par.022 | continuation; R28 par.022 visually replayed
\par\noindent
The equation $t^2-au^2=p$ gives immediately:
% U:p079.r28.eq.015 | continuation; R28 eq.015 visually replayed
\[
\begin{array}{llll}
\leg{-ap}{g}=1,& \leg{-ap}{g'}=1,& \leg{-ap}{g''}=1,& \text{etc.}\\[0.3em]
\leg{-ap}{h}=1,& \leg{-ap}{h'}=1,& \leg{-ap}{h''}=1,& \text{etc.}
\end{array}
\]
Comparing these relations with the preceding ones, one obtains:
% U:p079.r28.eq.016 | continuation; R28 eq.016 visually replayed
\[
\begin{array}{llll}
\leg{p}{g}=1,& \leg{p}{g'}=1,& \leg{p}{g''}=1,& \text{etc.}\\[0.3em]
\leg{p}{h}=-1,& \leg{p}{h'}=-1,& \leg{p}{h''}=-1,& \text{etc.}
\end{array}
\]
The application of the law of reciprocity will then give, $p$ being of the form $4n+1$:
% U:p079.r28.eq.017 | continuation; R28 eq.017 visually replayed
\[
\begin{array}{llll}
\leg{g}{p}=1,& \leg{g'}{p}=1,& \leg{g''}{p}=1,& \text{etc.}\\[0.3em]
\leg{h}{p}=-1,& \leg{h'}{p}=-1,& \leg{h''}{p}=-1,& \text{etc.}
\end{array}
\]
whence one concludes, by multiplying:
% U:p079.r28.eq.018a | continuation; R28 eq.018a visually replayed
\[
\leg{gg'g''\ldots{}\times hh'h''\ldots}{p}=\leg{t}{p}=\pm1,
\]
the upper or lower sign taking place according as the prime numbers $h$,
